\textbf{Note to the Judging Panel:} According to the initial
requirements of the internship program, these articles were expected to
be publicly posted on the community Facebook group. However, due to the
specific nature of the project (scraping and processing data from major
news websites) and to ensure data safety and security before the project
is 100\% finalized, our team decided \textbf{not to publish them
publicly on Facebook at this time}.

Instead, we have consolidated them as internal Blog posts within this
report. These represent the \textbf{valuable technical takeaways and
practical experiences that our entire team learned together} throughout
the process of building and optimizing the News RAG Pipeline system.

Below are 3 blog posts sharing experiences from the NewsRAG project:

\hypertarget{blog-1---overcoming-aws-lambda-timeouts-with-ecs-fargate}{%
\subsubsection{\texorpdfstring{\href{3.1-Blog1/}{Blog 1 - OVERCOMING AWS
LAMBDA TIMEOUTS WITH ECS
FARGATE}}{Blog 1 - OVERCOMING AWS LAMBDA TIMEOUTS WITH ECS FARGATE}}}

This article shares the practical experience of running the Crawler
system. Hitting Lambda's 15-minute limit prompted the team to flexibly
migrate to the ECS Fargate architecture, allowing the crawler to operate
stably for extended processing times.

\hypertarget{blog-2---system-cost-optimization-from-apache-kafka-to-amazon-sqs}{%
\subsubsection{\texorpdfstring{\href{3.2-Blog2/}{Blog 2 - SYSTEM COST
OPTIMIZATION: FROM APACHE KAFKA TO AMAZON
SQS}}{Blog 2 - SYSTEM COST OPTIMIZATION: FROM APACHE KAFKA TO AMAZON SQS}}}

This article analyzes the cost challenges when deploying on AWS. The
team bravely discarded the expensive Apache Kafka (which requires 24/7
server maintenance) in favor of Amazon SQS (Serverless), reducing the
queue component cost to nearly \$0.

\hypertarget{blog-3---realizing-a-closed-loop-rag-system-with-amazon-bedrock-and-aurora-pgvector}{%
\subsubsection{\texorpdfstring{\href{3.3-Blog3/}{Blog 3 - REALIZING A
CLOSED-LOOP RAG SYSTEM WITH AMAZON BEDROCK AND AURORA
PGVECTOR}}{Blog 3 - REALIZING A CLOSED-LOOP RAG SYSTEM WITH AMAZON BEDROCK AND AURORA PGVECTOR}}}

This article emphasizes security-minded and consistent architecture.
Instead of using a third-party Vector Database, the team successfully
integrated Amazon Bedrock and Aurora Serverless v2 + pgvector so the
entire RAG flow is processed securely within the AWS environment.
